% Paper-ready opening sections for the current XKV project.
% This file is intentionally standalone while the legacy EIYW draft is retained.
\documentclass[11pt]{article}

\usepackage[utf8]{inputenc}
\usepackage[T1]{fontenc}
\usepackage[margin=1in]{geometry}
\usepackage[hidelinks]{hyperref}
\usepackage{xspace}
\setlength{\emergencystretch}{2em}

\newcommand{\xkv}{\textsc{XKV}\xspace}
\newcommand{\lcfx}{\textsc{LCF-X}\xspace}
\newcommand{\ttoT}{\textsc{T2T}\xspace}
\newcommand{\sharer}{\textit{Sharer}\xspace}
\newcommand{\receiver}{\textit{Receiver}\xspace}

\title{\textbf{XKV}: Receiver-Native Latent Communication\\
Between Frozen Heterogeneous Language Models}
\author{Jiyao Liu\\
  \texttt{[affiliation and email]}}
\date{}

\begin{document}
\maketitle

\begin{abstract}
Systems of language-model agents usually communicate by generating text, even
when the sender and receiver are neural models whose internal states are
continuous. Although text is portable and interpretable, it introduces
autoregressive generation latency and restricts the exchanged information to a
discrete message selected without access to the receiver's internal state.
We introduce \xkv, a latent communication method for frozen, heterogeneous
language models that hold different evidence. A \sharer and a \receiver first
encode their private contexts independently. \xkv pools both key--value (KV)
caches with learned queries, aligns their layer geometries, and translates the
joint summaries into a compact layer memory. The receiver's raw cache positions
then query this memory to produce position-specific, per-head-scaled residuals
in the receiver's native KV space. Communication is therefore conditioned on
both what the sharer encoded and what the receiver already knows, while allowing
the two models to differ in family, depth, attention geometry, and tokenizer.
Only the lightweight translator is trained. On the complete $3\!\times\!3$
heterogeneous-model grid for ROPES, \xkv improves over the cross-context Latent
Cache Flow baseline by 4.6 exact-match and 4.2 F1 points on average, while using
77\% fewer trainable parameters and reducing end-to-end latency by 19\%.
Across seven evaluated HotpotQA-bridge model pairings, \xkv improves both exact
match and F1 in every pairing while reducing latency by approximately 20\%.
These results show that receiver-aware, receiver-native cache updates provide an
effective and efficient communication channel for models with complementary
contexts.
\end{abstract}

\section{Introduction}

Language models are increasingly deployed as components of larger systems
rather than as isolated predictors. Different agents may retrieve evidence,
invoke tools, propose solutions, or verify one another, and their collective
performance depends on how effectively information moves between them. Existing
multi-agent frameworks make natural language the common interface: agents
exchange dialogue turns, role-specific reports, or complete candidate answers
\cite{li2023camel,wu2023autogen,qian2023chatdev,hong2023metagpt,wang2024mixture}.
This interface is attractive because it works across model families and leaves a
human-readable trace. It is also an indirect neural communication protocol. The
sending model must decode its state into tokens one step at a time, and the
receiving model must tokenize and encode those tokens again before using them.

The limitations of this protocol are especially visible in the
\emph{cross-context} setting. Suppose a \sharer has one part of the evidence for
a question and a \receiver has another. Neither context is sufficient alone;
the receiver must integrate the missing evidence with its own context to answer.
A text message can communicate the missing facts, but its content is chosen from
the sharer's state alone. The sharer cannot observe how the receiver represented
its evidence, which ambiguity the receiver has resolved, or which internal
features would be most useful at each receiver layer. Moreover, text forces a
potentially rich internal representation through a finite vocabulary and adds a
sequential decode stage to the critical path.

Recent work has therefore explored communication through model activations.
Cache-to-Cache (C2C) translates one model's KV cache into another model's cache
space, but operates at corresponding token positions and consequently assumes
that both models process the same content \cite{fu2025cachetocache}. Latent Cache
Flow (LCF) compresses this transfer, and its cross-context extension, \lcfx,
pools the sharer's cache into position-free summaries before broadcasting them
over the receiver cache \cite{rossi2026latentcacheflow}. Other approaches pass
last-layer hidden states through a shared latent workspace
\cite{du2025interlat,zou2025latentmas}. These methods establish that text need
not be the only interface between language models. They also expose a central
design question: when the agents have different contexts, should the message be
a summary of the sender alone, or should it be computed relative to the state of
the model that must use it?

We pursue the latter view. We introduce \xkv, a dual-cache translator that
computes communication jointly from the sharer's and receiver's KV caches and
emits an update directly in the receiver's native cache geometry. After the two
frozen models independently prefill their private contexts, learned-query
attention reduces each layer and KV head to a small number of summaries. A
learned layer map reconciles different depths, and a bottleneck followed by
self-attention translates the concatenated summaries into a joint layer memory.
To restore information discarded by global pooling, each raw receiver cache
position acts as a query into this memory. A shared decoder then predicts
position-specific key and value residuals with deterministic per-head scaling.
The residual is applied once at the prompt boundary, after which the receiver
decodes normally. Thus, \xkv preserves the flexibility of cross-context
communication without requiring shared token positions, matched tokenizers, or
matched cache shapes.

This design has two practical consequences. First, the communication update is
\emph{receiver-aware}: it is a function of both caches rather than a
receiver-independent compression of the sharer's context. Second, the update is
\emph{receiver-native}: it is written into the representation that the frozen
receiver already uses for causal attention. Both base models remain unchanged;
training optimizes only the translator with answer-token supervision. The
translator's output projections are initialized to zero, making an untrained
module exactly reproduce receiver-only decoding and providing a clean
implementation and optimization boundary.

We evaluate cross-context communication on six complementary question-answering
tasks: HotpotQA-bridge, MuSiQue, QASC, 2WikiMultiHopQA, ROPES, and StrategyQA.
Their evidence is partitioned between the two agents, and the evaluation crosses
model families and parameter scales. The current complete ROPES grid shows that
\xkv improves over \lcfx by 4.6 exact-match and 4.2 F1 points on average, despite
using 77\% fewer trainable parameters; it also reduces end-to-end latency by
19\% and is approximately $8.9\times$ faster than bounded text communication.
On HotpotQA-bridge, \xkv improves both metrics in all seven evaluated pairings
and reduces average latency by roughly 20\%. Together, these results support
joint dual-cache translation as a compact interface for heterogeneous agents
with complementary information.

Our contributions are threefold:
\begin{itemize}
  \item We formulate cross-context latent communication as a joint translation
  problem over the sharer's and receiver's internal states, making the
  communicated update explicitly receiver-aware.
  \item We introduce \xkv, which combines symmetric direct cache pooling,
  cross-architecture layer alignment, attention-based translation, and a shared
  position decoder to generate receiver-native KV residuals while keeping both
  language models frozen.
  \item We provide a controlled evaluation across six split-evidence reasoning
  tasks and heterogeneous model pairs, comparing text communication and
  position-free latent cache fusion under a shared output and latency contract.
\end{itemize}

\section{Related Work}

\paragraph{Text-mediated multi-agent language models.}
Language is the default coordination channel for LLM agents. CAMEL studies
role-playing between communicative agents \cite{li2023camel}; AutoGen provides a
general framework for applications built from agent conversations
\cite{wu2023autogen}; and ChatDev and MetaGPT organize specialized agents into
language-mediated software-development workflows
\cite{qian2023chatdev,hong2023metagpt}. Mixture-of-Agents similarly passes the
textual outputs of one layer of models to the next \cite{wang2024mixture}.
These systems demonstrate the portability and interpretability of text, but
each handoff requires the sender to generate a discrete sequence and the
receiver to encode it. Our \ttoT baseline follows this conventional protocol.
\xkv instead communicates without an intermediate generation step, while
retaining compatibility across independently tokenized model families.

\paragraph{Continuous reasoning and latent agent communication.}
Continuous latent states have also been used to avoid committing intermediate
computation to language. Coconut feeds a model's last hidden state back as its
next input embedding, enabling reasoning in a continuous space
\cite{hao2024coconut}. Interlat transmits last hidden states between agents and
adds a learned latent compression process \cite{du2025interlat}, while LatentMAS
uses a shared latent working memory for training-free multi-agent collaboration
\cite{zou2025latentmas}. These approaches motivate continuous communication but
primarily operate on hidden-state sequences or a shared latent interface. \xkv
instead targets frozen heterogeneous models whose hidden dimensions, layer
counts, attention heads, head dimensions, and tokenizers may all differ. It
conditions the transfer on both agents' prefills and produces updates in the
receiver's own per-layer KV representation.

\paragraph{KV-cache transfer and reuse.}
The KV cache is a natural communication substrate because it is the persistent
state consumed during autoregressive decoding. C2C learns projections that fuse
a source cache into a target cache and gates the affected target layers
\cite{fu2025cachetocache}. Its position-wise fusion, however, relies on the two
models processing aligned versions of the same input. LCF reduces adapter cost
through a joint key--value bottleneck; \lcfx extends it to different contexts by
pooling the sharer's cache into position-free summaries before adding gated
receiver-cache residuals \cite{rossi2026latentcacheflow}. KVCOMM addresses a
different systems problem: it aligns cache offsets so overlapping text can be
reused across agents without retraining \cite{ye2025kvcomm}. \xkv is a learned
semantic transfer method rather than prefix reuse. Relative to \lcfx, it pools
both caches, mixes information across layers, and decodes the joint memory back
to position-specific receiver updates. This makes each update depend on the
receiver's encoded context rather than broadcasting a sharer-only summary.

\paragraph{Continuous prompts and context compression.}
Prefix tuning learns task-specific continuous vectors that a frozen language
model attends to as virtual tokens \cite{li2021prefixtuning}, and gist tokens
train a model to compress a prompt into a small reusable representation
\cite{mu2023gist}. Both lines show that continuous conditioning can be compact
and effective. Their goal differs from ours: a prefix is generally a static
task parameter, while gisting compresses context within a single model. \xkv
constructs an example-dependent signal from a separate sharer and the current
receiver state, translates across model geometries, and modifies the receiver's
existing KV cache rather than prepending embeddings or requiring the base model
to learn a compression behavior.

\bibliographystyle{plain}
\bibliography{references}

\end{document}
